Extend $AE$ and $CD$ to meet at $F$, completing the parallelogram whose other vertices are $A,B,C$. Opposite sides give $AF=CF=10$, hence $EF=DF=2$. In triangle $DEF$, the cosine rule gives $\cos\angle F=\dfrac{(4+4-9)}{8}=-\dfrac{1}{8}$. To connect this with the cube side $s$, put the cutting plane in the form $z=\alpha x+\beta y+\gamma$ across two adjacent side-face widths. The two full parallelogram sides have lengths $s\sqrt{1+\alpha^2}=s\sqrt{1+\beta^2}=10$, so $|\alpha|=|\beta|$. Their obtuse angle has cosine $-\dfrac{\alpha^2}{(1+\alpha^2)}=-\dfrac{1}{8}$, giving $\alpha^2=\dfrac{1}{7}$. Thus $s^2=\dfrac{100}{(1+\dfrac{1}{7})}=\dfrac{175}{2}$ and the cube surface area is $6s^2=\boxed{525}$.
\begin{asy}
import graph;import geometry;usepackage("amsmath");size(7cm);
pair F=(0,0),A=(10,0),C=(-1.25,10*sqrt(63)/8),B=A+C,D=C/5,E=A/5;draw(A--B--C--D--E--cycle);draw(D--F--E,dashed);label("$A$",A,(1,-1));label("$B$",B,(1,1));label("$C$",C,(-1,1));label("$D$",D,(-1,0));label("$E$",E,(0,-1));label("$F$",F,(-1,-1));label("$10$",(A+B)/2,(1,0));label("$10$",(B+C)/2,(0,1));label("$8$",(C+D)/2,(-1,0));label("$8$",(A+E)/2,(0,-1));label("$3$",(D+E)/2,(1,1));
\end{asy}

